% !TEX TS-program = pdflatex
% CV ciblé — Ingénieur Logiciel IA — AZ Information Technology (plateforme EdTech IA)
\documentclass[a4paper,9pt]{article}
\usepackage[utf8]{inputenc}
\usepackage[T1]{fontenc}
\usepackage[scaled=0.90]{helvet}
\renewcommand{\familydefault}{\sfdefault}
\usepackage[left=0.8cm,right=0.8cm,top=0.3cm,bottom=0.25cm]{geometry}
\usepackage{enumitem}
\usepackage{titlesec}
\usepackage{xcolor}
\usepackage[hidelinks]{hyperref}
\definecolor{navy}{RGB}{23,54,93}
\pagestyle{empty}
\setlength{\parindent}{0pt}
\setlength{\parskip}{0pt}
\linespread{0.82}
\hypersetup{pdftitle={CV - Baha Eddine Belhaj Mouldi - Ingenieur Logiciel IA - AZ Information Technology},pdfauthor={Baha Eddine Belhaj Mouldi},pdfsubject={IA, LLM, RAG, FastAPI, PostgreSQL, Azure, Backend, EdTech}}
\titleformat{\section}{\normalsize\bfseries\color{navy}}{}{0pt}{\MakeUppercase}[{\vspace{1pt}\color{navy}\hrule height 0.6pt}]
\titlespacing*{\section}{0pt}{1.5pt}{0.8pt}
\setlist[itemize]{leftmargin=1.0em,label=\textbullet,itemsep=0pt,topsep=0pt,parsep=0pt}
\newcommand{\cventry}[4]{\textbf{#1} \hfill \textbf{#4}\\[0.2pt]#2{\ifx\relax#3\relax\else, #3\fi}\par\vspace{0.3pt}}
\begin{document}
\begin{center}
{\LARGE\bfseries\color{navy} Baha Eddine Belhaj Mouldi}\\[2pt]
{\normalsize\color{navy} Ingénieur Logiciel IA \textbar\ Backend Python \textbar\ RAG/LLM \textbar\ Azure AI \textbar\ FastAPI}\\[3pt]
{\small Tunis, Tunisie \quad|\quad +216 55 128 982 \quad|\quad \href{mailto:bahaeddine.belhajmouldi@gmail.com}{bahaeddine.belhajmouldi@gmail.com}}\\[1pt]
{\small \href{https://www.linkedin.com/in/bahamouldi}{linkedin.com/in/bahamouldi} \quad|\quad \href{https://github.com/bahamouldi}{github.com/bahamouldi} \quad|\quad \href{https://bahamouldi.github.io/portfolio/}{bahamouldi.github.io/portfolio}}
\end{center}
\vspace{1pt}
\section{Profil}
Ingénieur diplômé en génie informatique (ENIT, mention Très Bien) spécialisé dans la conception de services backend IA : pipelines RAG, intégration de modèles de langage, architectures multi-agents et APIs FastAPI. Expérience concrète et déjà en production chez Beehive Entreprises (stage PFE et un an de freelance) : chatbot RAG sur 900K enregistrements, orchestration LLM (Gemini, Azure AI Foundry), moteurs d'évaluation de la qualité de génération (LangSmith), déploiement Kubernetes sécurisé et scalable. Pratique concrète d'Azure AI Agent Service et du Microsoft Agent Framework SDK (protocoles Agent-to-Agent, agents collaboratifs) au travers de projets personnels. Base solide en gestion de données (PostgreSQL) et en sécurité applicative (OWASP, gestion des correctifs). Curieux des sujets EdTech et de l'apprentissage des langues --- pratique personnelle du chinois. Bilingue français/anglais, ouverture à la mobilité internationale.
\section{Compétences clés}
\textbf{Backend \& APIs} : Python, FastAPI, REST APIs, architecture microservices, async, gestion des erreurs et observabilité\\[0.3pt]
\textbf{Intégration IA \& LLM} : RAG (retrieval, injection de contexte, génération, évaluation), LangChain, LangGraph, Gemini, Azure AI Foundry, Azure AI Agent Service, prompt engineering et versionning\\[0.3pt]
\textbf{Moteurs d'évaluation \& rétroaction} : LangSmith (tracing, évaluation, monitoring qualité), métriques de latence/precision, dashboards de suivi (Grafana, Streamlit)\\[0.3pt]
\textbf{Agents \& orchestration} : architectures multi-agents, protocole Agent-to-Agent (A2A), Microsoft Agent Framework SDK, MCP, orchestration de workflows\\[0.3pt]
\textbf{Bases de données \& cloud} : PostgreSQL, SQL, MongoDB, Azure (AI Foundry, AI Search, Functions), Docker, Kubernetes, CI/CD (Jenkins, ArgoCD, GitHub Actions)\\[0.3pt]
\textbf{Machine Learning} : Scikit-learn, Pandas, SHAP, NLP, classification de texte, analyse de sentiments\\[0.3pt]
\textbf{Sécurité \& scalabilité} : OWASP Top 10/API/LLM, gestion des correctifs (CVE), conformité (ISO 27001, GDPR), déploiement zéro downtime\\[0.3pt]
\textbf{Outils} : Git, documentation technique (Swagger/OpenAPI), revue de code, Bash, Linux
\section{Expérience professionnelle}
\cventry{Ingénieur Logiciel IA --- Stagiaire PFE (mention Très Bien)}{Beehive Entreprises}{Tunisie}{02/2026 -- 06/2026}
\begin{itemize}
  \item Conception de services backend IA en production : chatbot RAG sur 900K enregistrements (LangChain, Gemini, Azure AI Foundry), architectures multi-agents, APIs FastAPI déployées sur Kubernetes.
  \item Moteur d'évaluation et de rétroaction : LangSmith (tracing, évaluation de la qualité de génération), dashboards Prometheus/Grafana, alerting temps réel sur les performances des agents.
  \item Sécurité et scalabilité : conformité continue sur 7 référentiels (OWASP, PCI DSS, GDPR, NIST, ISO 27001), pipeline CI/CD (Jenkins + ArgoCD), déploiement zéro downtime.
\end{itemize}
\cventry{Ingénieur IA \& Backend, Freelance (1 an)}{Beehive Entreprises}{Tunisie}{01/2025 -- 12/2025}
\begin{itemize}
  \item Développement de services backend IA pour applications clients : FastAPI, intégrations LLM, pipelines RAG, orchestration multi-agents, gestion PostgreSQL.
  \item Revue de code et analyse de vulnérabilités sur 8 applications web et API ; collaboration transverse avec les équipes métier.
\end{itemize}
\cventry{Stagiaire Cybersécurité --- Détection menaces SAP}{Streamlink}{Tunisie}{07/2025 -- 08/2025}
\begin{itemize}
  \item Pipeline de détection automatisé (Python, corrélation ELK, orchestration N8N) --- workflows testables et auditables.
\end{itemize}
\cventry{Chef de projet technique, Indépendant --- Équipe de 8 personnes}{VMate}{Tunisie}{01/2024 -- 06/2024}
\begin{itemize}
  \item Pilotage d'une équipe pluridisciplinaire de 8 personnes : planification, arbitrages techniques, reporting.
\end{itemize}
\section{Projets clés}
\cventry{Chatbot RAG --- Dashboard Assurance Revenus (Tunisie Telecom)}{Python, LangChain, Gemini, Azure AI Foundry, PostgreSQL}{}{2026}
\begin{itemize}
  \item Pipeline RAG en production sur ~900K enregistrements : retrieval, injection de contexte, génération. Interrogation en langage naturel pour utilisateurs métier, monitoring LangSmith/Grafana.
\end{itemize}
\cventry{Agents Distants avec Protocole Azure AI Agent Service (A2A)}{Python, Azure AI Foundry, Azure AI Agent Service}{}{2026}
\begin{itemize}
  \item Agents collaboratifs communiquant via le protocole Agent-to-Agent sur Azure AI Foundry : génération de contenu et enchaînement de tâches entre agents spécialisés.
\end{itemize}
\cventry{Solution Multi-Agent --- Microsoft Agent Framework}{Python, Microsoft Agent Framework SDK, Azure AI Foundry}{}{2026}
\begin{itemize}
  \item Pipeline séquentiel de 3 agents (résumé, classification, recommandation d'action) pour le traitement de feedback client --- orchestration testable et modulaire.
\end{itemize}
\cventry{Système de Recommandation d'Hôtels RAG}{Python, LangChain, LLM, Azure AI Foundry}{}{2025}
\begin{itemize}
  \item Recherche sémantique sur avis clients, recommandations contextuelles via LLM, orchestration multi-agents.
\end{itemize}
\cventry{BeeWAF --- Sécurité Applicative \& Scalabilité (PFE, 86/100)}{FastAPI, Python, ML, Docker, Kubernetes, ELK}{}{2026}
\begin{itemize}
  \item Backend FastAPI sécurisé et scalable en production : 27 modules, 3 modèles ML, CI/CD, conformité continue, zéro downtime.
\end{itemize}
\section{Formation}
\cventry{Diplôme d'Ingénieur en Génie Informatique --- Mention Très Bien}{École Nationale d'Ingénieurs de Tunis (ENIT)}{Tunisie}{09/2023 -- 06/2026}
\begin{itemize}
  \item Diplômé le 25/06/2026. Classé 65e sur 600+ au concours national d'entrée. PFE : architecture multi-agents et RAG (86/100).
\end{itemize}
\cventry{Classes préparatoires Physique-Technologie}{Institut Préparatoire aux Études d'Ingénieurs, El Manar}{Tunisie}{09/2021 -- 06/2023}
\section{Certifications}
Fondements du Deep Learning --- NVIDIA (2025) ; IA Adversariale --- NVIDIA (2025) ; Machine Learning --- NVIDIA (2024) ; Réseaux (CCNA) --- Cisco (2025) ; Git \& GitHub --- 365 Data Science (2024).
\section{Langues}
Français (Courant) \quad|\quad Anglais (B2) \quad|\quad Arabe (Natif) \quad|\quad Chinois (Notions)
\end{document}
